\ELhold

\ELsection{Course Outline}{Course Outline}

\begin{ELunit}

\ELfigure{m05-course-outline}{}

\end{ELunit}

\ELhold

\ELsection{Static Magnetic Fields}{Static Magnetic Fields}

\begin{ELunit}

\begin{itemize}
\tightlist
\item
  When a small test charge is placed in an electric field \ELim{\vect{E}}, the following force acts on it
\end{itemize}

\ELdisp{\vect{F}_e\ELbr =q\vect{E}\qquad\ELbr (\mathrm{N})}

\begin{itemize}
\tightlist
\item
  When this test charge moves in a magnetic field, another force acts on it, equal to
\end{itemize}

\ELdisp{\vect{F}_m\ELbr =q\vect{u}\times\vect{B}\qquad\ELbr (\mathrm{N})}

\end{ELunit}

\begin{ELjoin}

\begin{itemize}
\tightlist
\item
  \ELim{\vect{B}}: the magnetic flux density, in webers per square metre or teslas
\item
  So the total electromagnetic force on the charge \ELim{q} is
\end{itemize}

\begin{ELimportant}{Lorentz's force equation}

\ELdisp{\vect{F}\ELbr =q(\vect{E}+\vect{u}\times\vect{B})\qquad\ELbr (\mathrm{N})}

\end{ELimportant}

\end{ELjoin}

\begin{ELunit}

\ELfigure{m05-map}{}

\begin{itemize}
\tightlist
\item
  What are the fundamental postulates of magnetostatics, from which the other relations and laws of the field can be derived?
\item
  How is \textbf{Ampère's circuital law} derived from these postulates?
\item
  How can the magnetic flux density of a current be found with Ampère's circuital law?
\end{itemize}

\end{ELunit}

\ELhold

\ELsection{Fundamental Postulates of Magnetostatics in Free Space}{Fundamental Postulates of Magnetostatics in Free Space}

\begin{ELjoin}

\begin{itemize}
\tightlist
\item
  The two fundamental postulates of magnetostatics in free space
\end{itemize}

\begin{ELimportant}{}

\ELdisp{\nabla\cdot\vect{B}\ELbr =0}
\ELdisp{\nabla\times\vect{B}\ELbr =\mu_0\vect{J}}

\end{ELimportant}

\end{ELjoin}

\begin{ELunit}

\begin{itemize}
\tightlist
\item
  \ELim{\mu_0=4\pi\times10^{-7}\,\mathrm{(H/m)}}: the permeability of free space
\item
  \ELim{\vect{J}}: the volume current density \ELim{\mathrm{(A/m^2)}}
\item
  Since the divergence of the curl of any vector field is zero, it follows that
\end{itemize}

\ELdisp{\nabla\cdot\vect{J}\ELbr =0}

\end{ELunit}

\begin{ELjoin}

\begin{itemize}
\tightlist
\item
  There is no magnetic analogue of the electric charge density.
\item
  The integral form of the first relation \ELim{\ELto} integrate over an arbitrary volume \ELim{V} and apply the divergence theorem
\end{itemize}

\begin{ELimportant}{Law of conservation of magnetic flux}

\ELdisp{\oint_S\vect{B}\cdot d\vect{s}\ELbr =0}

\end{ELimportant}

\end{ELjoin}

\begin{ELunit}

\begin{itemize}
\tightlist
\item
  This equation is called the \textbf{law of conservation of magnetic flux}.

  \begin{itemize}
  \tightlist
  \item
    There are no magnetic flow sources, and the magnetic flux lines always close upon themselves.
  \end{itemize}
\item
  The integral form of the second relation \ELim{\ELto} integrate over an arbitrary open surface \ELim{S} and apply Stokes's theorem
\end{itemize}

\ELdisp{\int_S(\nabla\times\vect{B})\cdot d\vect{s}\ELbr =\mu_0\int_S\vect{J}\cdot d\vect{s}}

\end{ELunit}

\begin{ELimportant}{Ampère's circuital law}

\ELdisp{\oint_C\vect{B}\cdot\dif\vect{\ell}\ELbr =\mu_0I}

\end{ELimportant}

\begin{ELunit}

\begin{itemize}
\tightlist
\item
  \ELim{C}: the contour bounding the surface \ELim{S} (the contour \ELim{C} and the current \ELim{I} follow the right-hand rule)
\item
  \ELim{I}: the current passing through the surface \ELim{S}
\item
  This equation is a form of \textbf{Ampère's circuital law}.

  \begin{itemize}
  \tightlist
  \item
    The circulation of the magnetic flux density in free space around any closed path is equal to \ELim{\mu_0} times the total current flowing through the surface bounded by the path.
  \end{itemize}
\item
  Ampère's circuital law is useful for finding the magnetic flux density caused by a current \ELim{I} when there is a closed path \ELim{C} around the current such that \ELim{\vect{B}} is constant on the path.
\end{itemize}

\end{ELunit}

\ELnextnum{5-1}

\begin{ELexample}{}

An infinitely long, straight conductor with a circular cross-section of radius \ELim{b} carries a current \ELim{I}. Determine the magnetic flux density both inside and outside the conductor.

\ELfigure{m05-ex1}{}

\begin{ELsolution}

\begin{itemize}
\tightlist
\item
  If the conductor lies along the \ELim{z}-axis, \ELim{\vect{B}} is in the \ELim{\phi} direction and its magnitude is constant along any circular path about the \ELim{z}-axis.

  \begin{itemize}
  \tightlist
  \item
    Inside the conductor
  \end{itemize}
\end{itemize}

\ELdisp{\left.\begin{aligned}&\vect{B}_1=\uvec{\phi}B_{\phi1},\qquad\dif\vect{\ell}=\uvec{\phi}r_1\,d\phi\\&\oint_{C_1}\vect{B}_1\cdot\dif\vect{\ell}=\int_0^{2\pi}B_{\phi1}r_1\,d\phi=2\pi r_1B_{\phi1}\\&I_1=\frac{\pi r_1^2}{\pi b^2}I=\left(\frac{r_1}{b}\right)^2I\end{aligned}\right\}\quad\ELbr \vect{B}_1\ELbr =\uvec{\phi}B_{\phi1}\ELbr =\uvec{\phi}\frac{\mu_0r_1I}{2\pi b^2}}

\begin{itemize}
\tightlist
\item
  Outside the conductor
\end{itemize}

\ELdisp{\vect{B}_2\ELbr =\uvec{\phi}B_{\phi2},\qquad\ELbr \dif\vect{\ell}\ELbr =\uvec{\phi}r_2\,d\phi}
\ELdisp{\oint_{C_2}\vect{B}_2\cdot\dif\vect{\ell}\ELbr =2\pi r_2B_{\phi2}}
\ELdisp{\vect{B}_2\ELbr =\uvec{\phi}B_{\phi2}\ELbr =\uvec{\phi}\frac{\mu_0I}{2\pi r_2}}

\end{ELsolution}

\end{ELexample}

\ELnextnum{5-2}

\begin{ELexample}{}

Determine the magnetic flux density inside a closely wound toroidal coil with an air core, having \ELim{N} turns of wire carrying a current \ELim{I}.

\ELfigure{m05-ex2}{}

\begin{ELsolution}

\begin{itemize}
\tightlist
\item
  Cylindrical symmetry ensures that \ELim{\vect{B}} has only a \ELim{\phi} component and that its magnitude is constant along any circular path about the axis of the coil.
\end{itemize}

\ELdisp{\oint\vect{B}\cdot\dif\vect{\ell}\ELbr =2\pi rB_\phi\ELbr =\mu_0NI}
\ELdisp{\vect{B}\ELbr =\uvec{\phi}B_\phi\ELbr =\uvec{\phi}\frac{\mu_0NI}{2\pi r},\qquad\ELbr (b-a)<r<(b+a)}

\end{ELsolution}

\end{ELexample}

\ELnextnum{5-3}

\begin{ELexample}{}

Determine the magnetic flux density inside an infinitely long solenoid with an air core, having \ELim{n} turns of wire per unit length carrying a current \ELim{I}.

\ELfigure{m05-ex3}{}

\begin{ELsolution}

\begin{itemize}
\tightlist
\item
  Outside the solenoid the magnetic field is zero.
\item
  By the symmetry of the structure, \ELim{\vect{B}} is parallel to the axis of the solenoid and its magnitude is constant along the solenoid.
\end{itemize}

\ELdisp{BL\ELbr =\mu_0nLI}
\ELdisp{B\ELbr =\mu_0nI}

\end{ELsolution}

\end{ELexample}

\ELhold

\ELsection{Vector Magnetic Potential}{Vector Magnetic Potential}

\begin{ELunit}

\ELfigure{m05-map-vecpot}{}

\begin{itemize}
\tightlist
\item
  What is the vector magnetic potential, where is it used, and how is it calculated?
\item
  What is the Biot--Savart law, and how is it used to calculate the magnetic flux density?
\end{itemize}

\ELdisp{\nabla\cdot\vect{B}\ELbr =0\quad\ELbr \ELbr \Longrightarrow\quad\ELbr \boxed{\vect{B}=\nabla\times\vect{A}\qquad(\mathrm{T})}}

\begin{itemize}
\tightlist
\item
  \ELim{\vect{A}} is called the \textbf{vector magnetic potential}. Its unit is the weber per metre.

  \begin{itemize}
  \tightlist
  \item
    If \ELim{\vect{A}} can be found for a current distribution, \ELim{\vect{B}} follows easily (exactly as with \ELim{V} and \ELim{\vect{E}}).
  \end{itemize}
\item
  A vector is defined by its curl and its divergence. To specify the divergence of \ELim{\vect{A}} we proceed as follows
\end{itemize}

\ELdisp{\nabla\times\vect{B}\ELbr =\mu_0\vect{J}\quad\ELbr \ELbr \Longrightarrow\quad\ELbr \nabla\times\nabla\times\vect{A}\ELbr =\mu_0\vect{J}}

\begin{itemize}
\tightlist
\item
  We define the \textbf{Laplacian of a vector quantity} as
\end{itemize}

\ELdisp{\nabla^2\vect{A}\ELbr =\nabla(\nabla\cdot\vect{A})-\nabla\times\nabla\times\vect{A}}

\begin{itemize}
\tightlist
\item
  For example, the Laplacian of \ELim{\vect{A}} in Cartesian coordinates is
\end{itemize}

\ELdisp{\nabla^2\vect{A}\ELbr =\uvec{x}\nabla^2A_x+\uvec{y}\nabla^2A_y+\uvec{z}\nabla^2A_z}

\begin{itemize}
\tightlist
\item
  Therefore
\end{itemize}

\ELdisp{\nabla(\nabla\cdot\vect{A})-\nabla^2\vect{A}\ELbr =\mu_0\vect{J}}

\begin{itemize}
\tightlist
\item
  To simplify this relation we choose
\end{itemize}

\ELdisp{\nabla\cdot\vect{A}\ELbr =0\quad\ELbr \ELbr \Longrightarrow\quad\ELbr \boxed{\nabla^2\vect{A}=-\mu_0\vect{J}}}

\begin{itemize}
\tightlist
\item
  This last relation is the \textbf{vector Poisson's equation}.
\item
  In Cartesian coordinates
\end{itemize}

\ELdisp{\nabla^2A_x\ELbr =-\mu_0J_x,\qquad\ELbr \nabla^2A_y\ELbr =-\mu_0J_y,\qquad\ELbr \nabla^2A_z\ELbr =-\mu_0J_z}

\begin{itemize}
\tightlist
\item
  Each of these three equations is mathematically the same as Poisson's equation in electrostatics.
\item
  Earlier we had
\end{itemize}

\ELdisp{\nabla^2V\ELbr =-\frac{\rho}{\epsilon_0}\quad\ELbr \ELbr \Longrightarrow\quad\ELbr  V\ELbr =\frac{1}{4\pi\epsilon_0}\int_{V'}\frac{\rho}{\lvert\vect{R}-\vect{R}'\rvert}\,dv'}

\end{ELunit}

\begin{ELjoin}

\begin{itemize}
\tightlist
\item
  In the same way we obtain
\end{itemize}

\begin{ELimportant}{}

\ELdisp{\vect{A}\ELbr =\frac{\mu_0}{4\pi}\int_{V'}\frac{\vect{J}}{\lvert\vect{R}-\vect{R}'\rvert}\,dv'}

\end{ELimportant}

\end{ELjoin}

\begin{ELjoin}

\begin{itemize}
\tightlist
\item
  So \ELim{\vect{A}} can be found from the current density with this relation, and \ELim{\vect{B}} then follows by taking its curl.
\item
  The physical meaning of the vector magnetic potential
\end{itemize}

\ELdisp{\Phi\ELbr =\int_S\vect{B}\cdot d\vect{s}}

\begin{itemize}
\tightlist
\item
  \ELim{\Phi}: the magnetic flux, in webers
\end{itemize}

\begin{ELimportant}{}

\ELdisp{\Phi\ELbr =\int_S(\nabla\times\vect{A})\cdot d\vect{s}\ELbr =\oint_C\vect{A}\cdot\dif\vect{\ell}\qquad\ELbr (\mathrm{Wb})}

\end{ELimportant}

\end{ELjoin}

\begin{ELunit}

\begin{itemize}
\tightlist
\item
  The line integral of \ELim{\vect{A}} around any closed path \ELim{C} equals the total magnetic flux passing through the surface bounded by the path.
\end{itemize}

\end{ELunit}

\ELhold

\ELsection{The Biot--Savart Law}{The Biot–Savart Law}

\begin{ELunit}

\ELfigure{m05-map-biot}{}

\begin{itemize}
\tightlist
\item
  For a thin wire carrying a current \ELim{I}, with cross-section \ELim{S}, we have
\end{itemize}

\ELdisp{\vect{J}\,dv'\ELbr =JS\,d\ell'\ELbr =I\,\dif\vect{\ell}'\quad\ELbr \ELbr \Longrightarrow\quad\ELbr \vect{A}\ELbr =\frac{\mu_0I}{4\pi}\oint_{C'}\frac{\dif\vect{\ell}'}{\lvert\vect{R}-\vect{R}'\rvert}}

\begin{itemize}
\tightlist
\item
  It can be shown that
\end{itemize}

\ELdisp{\vect{B}\ELbr =\nabla\times\vect{A}\ELbr =\nabla\times\left[\frac{\mu_0I}{4\pi}\oint_{C'}\frac{\dif\vect{\ell}'}{\lvert\vect{R}-\vect{R}'\rvert}\right]}

\end{ELunit}

\begin{ELimportant}{Biot--Savart law}

\ELdisp{\vect{B}\ELbr =\frac{\mu_0I}{4\pi}\oint_{C'}\frac{\dif\vect{\ell}'\times(\vect{R}-\vect{R}')}{\lvert\vect{R}-\vect{R}'\rvert^3}}

\end{ELimportant}

\begin{ELjoin}

\begin{itemize}
\tightlist
\item
  This equation is called the \textbf{Biot--Savart law}.

  \begin{itemize}
  \tightlist
  \item
    In general the Biot--Savart law is more difficult to use than Ampère's circuital law. But when no closed path can be found on which \ELim{\vect{B}} is constant, Ampère's circuital law cannot be used to determine \ELim{\vect{B}}.
  \end{itemize}
\end{itemize}

\ELnextnum{5-4}

\begin{ELexample}{}

A direct current \ELim{I} flows in a straight wire of length \ELim{2L}. Find the magnetic flux density at a point at a distance \ELim{r} from the wire in the bisecting plane.

\ELfigure{m05-ex4}{}

\begin{ELsolution}

\begin{itemize}
\tightlist
\item
  Currents exist only in closed circuits. The wire in this example must therefore be part of a current-carrying loop. Since we know nothing about the rest of the circuit, Ampère's circuital law cannot be used.
\item
  \textbf{First method}
\end{itemize}

\ELdisp{\left.\begin{aligned}&\vect{A}=\frac{\mu_0I}{4\pi}\oint_{C'}\frac{\dif\vect{\ell}'}{\lvert\vect{R}-\vect{R}'\rvert}\\&\dif\vect{\ell}'=dz'\,\uvec{z}\\&\vect{R}=r\uvec{r}\\&\vect{R}'=z'\uvec{z}\end{aligned}\right\}\quad\ELbr \begin{aligned}\vect{A}&=\uvec{z}\frac{\mu_0I}{4\pi}\int_{-L}^{L}\frac{dz'}{\sqrt{z'^2+r^2}}\\&=\uvec{z}\frac{\mu_0I}{4\pi}\ln\frac{\sqrt{L^2+r^2}+L}{\sqrt{L^2+r^2}-L}\end{aligned}}
\ELdisp{\vect{B}\ELbr =\nabla\times\vect{A}\ELbr =\nabla\times(\uvec{z}A_z)\ELbr =\uvec{r}\frac1r\frac{\partial A_z}{\partial\phi}-\uvec{\phi}\frac{\partial A_z}{\partial r}}
\ELdisp{\begin{aligned}\vect{B}&=-\uvec{\phi}\frac{\partial}{\partial r}\left[\frac{\mu_0I}{4\pi}\ln\frac{\sqrt{L^2+r^2}+L}{\sqrt{L^2+r^2}-L}\right]\\&=\uvec{\phi}\frac{\mu_0IL}{2\pi r\sqrt{L^2+r^2}}\end{aligned}}

\begin{itemize}
\tightlist
\item
  \textbf{Second method}
\end{itemize}

\ELdisp{\left.\begin{aligned}&\vect{B}=\frac{\mu_0I}{4\pi}\oint_{C'}\frac{\dif\vect{\ell}'\times(\vect{R}-\vect{R}')}{\lvert\vect{R}-\vect{R}'\rvert^3}\\&\vect{R}-\vect{R}'=r\uvec{r}-z'\uvec{z}\\&\dif\vect{\ell}'\times(\vect{R}-\vect{R}')=dz'\uvec{z}\times(r\uvec{r}-z'\uvec{z})=r\,dz'\,\uvec{\phi}\end{aligned}\right\}\quad\ELbr \begin{aligned}\vect{B}=\int d\vect{B}&=\uvec{\phi}\frac{\mu_0I}{4\pi}\int_{-L}^{L}\frac{r\,dz'}{(z'^2+r^2)^{3/2}}\\&=\uvec{\phi}\frac{\mu_0IL}{2\pi r\sqrt{L^2+r^2}}\end{aligned}}

\end{ELsolution}

\end{ELexample}

\end{ELjoin}

\ELnextnum{5-5}

\begin{ELexample}{}

Find the magnetic flux density at the centre of a square loop of side \ELim{w} carrying a current \ELim{I}.

\ELfigure{m05-ex5}{}

\begin{ELsolution}

\begin{itemize}
\tightlist
\item
  Assume the loop lies in the \ELim{xy}-plane. \ELim{\vect{B}} at the centre of the loop is four times the magnetic flux density caused by one side of length \ELim{w}. Therefore
\item
  The magnitude of the magnetic flux density of a straight wire of length \ELim{2L} at a distance \ELim{r} from it
\end{itemize}

\ELdisp{B\ELbr =\frac{\mu_0IL}{2\pi r\sqrt{L^2+r^2}}\qquad\ELbr  L\ELbr =\frac w2,\quad\ELbr  r\ELbr =\frac w2}
\ELdisp{\vect{B}\ELbr =\uvec{z}\frac{\mu_0I}{\sqrt2\pi w}\times4\ELbr =\uvec{z}\frac{2\sqrt2\mu_0I}{\pi w}}

\end{ELsolution}

\end{ELexample}

\ELnextnum{5-6}

\begin{ELexample}{}

Find the magnetic flux density on the axis of a circular loop of radius \ELim{b} carrying a current \ELim{I}.

\ELfigure{m05-ex6}{}

\begin{ELsolution}

\ELdisp{\left.\begin{aligned}&\vect{B}=\frac{\mu_0I}{4\pi}\oint_{C'}\frac{\dif\vect{\ell}'\times(\vect{R}-\vect{R}')}{\lvert\vect{R}-\vect{R}'\rvert^3}\\&\dif\vect{\ell}'=b\,d\phi'\,\uvec{\phi'},\quad\vect{R}=z\uvec{z},\quad\vect{R}'=b\uvec{r'}\\&\vect{R}-\vect{R}'=z\uvec{z}-b\uvec{r'}\\&\begin{aligned}\dif\vect{\ell}'\times(\vect{R}-\vect{R}')&=b\,d\phi'\,\uvec{\phi'}\times(z\uvec{z}-b\uvec{r'})\\&=bz\,d\phi'\,\uvec{r'}+b^2d\phi'\,\uvec{z}\end{aligned}\end{aligned}\right\}\quad\ELbr \vect{B}\ELbr =\frac{\mu_0I}{4\pi}\int_0^{2\pi}\uvec{z}\frac{b^2\,d\phi'}{(z^2+b^2)^{3/2}}}

\begin{itemize}
\tightlist
\item
  It is evident from the geometry that the \ELim{r}-component of the magnetic flux density is zero.
\end{itemize}

\begin{ELimportant}{}

\ELdisp{\vect{B}\ELbr =\uvec{z}\frac{\mu_0Ib^2}{2(z^2+b^2)^{3/2}}\qquad\ELbr (\mathrm{T})}

\end{ELimportant}

\end{ELsolution}

\end{ELexample}

\ELnextnum{5-7}

\begin{ELexample}{}

Find the magnetic flux density at a distant point of a magnetic dipole of radius \ELim{b} carrying a current \ELim{I}.

\ELfigure{m05-ex7}{}

\begin{ELsolution}

\ELdisp{\vect{A}\ELbr =\frac{\mu_0I}{4\pi}\oint_{C'}\frac{\dif\vect{\ell}'}{\lvert\vect{R}-\vect{R}'\rvert}}

\begin{itemize}
\tightlist
\item
  Because of the symmetry of the geometry, the magnetic field is independent of the angle \ELim{\phi}. For simplicity, the field point is therefore placed in the \ELim{yz}-plane.
\end{itemize}

\ELdisp{\dif\vect{\ell}'\ELbr =b\,d\phi'\,\uvec{\phi'}\ELbr =b\,d\phi'\left(-\uvec{x}\sin\phi'+\uvec{y}\cos\phi'\right)}
\ELdisp{\vect{A}\ELbr =-\uvec{x}\frac{\mu_0I}{4\pi}\int_0^{2\pi}\frac{b\sin\phi'}{\lvert\vect{R}-\vect{R}'\rvert}\,d\phi'+\uvec{y}\frac{\mu_0I}{4\pi}\int_0^{2\pi}\frac{b\cos\phi'}{\lvert\vect{R}-\vect{R}'\rvert}\,d\phi'}

\begin{itemize}
\tightlist
\item
  It is evident from the geometry that the \ELim{y}-component of \ELim{\vect{A}} is zero.
\end{itemize}

\ELfigure{m05-ex7-top}{}

\ELdisp{\begin{aligned}\vect{R}=R\uvec{R}&=R\left(\sin\theta\cos(\pi/2)\uvec{x}+\sin\theta\sin(\pi/2)\uvec{y}+\cos\theta\,\uvec{z}\right)\\&=R\left(\sin\theta\,\uvec{y}+\cos\theta\,\uvec{z}\right)\end{aligned}}
\ELdisp{\vect{R}'\ELbr =b\uvec{r'}\ELbr =b\left(\cos\phi'\,\uvec{x}+\sin\phi'\,\uvec{y}\right)}
\ELdisp{\vect{R}-\vect{R}'\ELbr =-b\cos\phi'\,\uvec{x}+(R\sin\theta-b\sin\phi')\uvec{y}+R\cos\theta\,\uvec{z}}
\ELdisp{\begin{aligned}\lvert\vect{R}-\vect{R}'\rvert^2&=R^2+b^2-2bR\sin\theta\sin\phi'\\&=R^2\left(1+\frac{b^2}{R^2}-\frac{2b}{R}\sin\theta\sin\phi'\right)\end{aligned}}
\ELdisp{\begin{aligned}\frac{1}{\lvert\vect{R}-\vect{R}'\rvert}&=\frac1R\left(1+\frac{b^2}{R^2}-\frac{2b}{R}\sin\theta\sin\phi'\right)^{-1/2}\\&\cong\frac1R\left(1-\frac{2b}{R}\sin\theta\sin\phi'\right)^{-1/2}\cong\frac1R\left(1+\frac bR\sin\theta\sin\phi'\right)\end{aligned}}
\ELdisp{\begin{aligned}\vect{A}&=-\uvec{x}\frac{\mu_0Ib}{4\pi R}\int_0^{2\pi}\left(1+\frac bR\sin\theta\sin\phi'\right)\sin\phi'\,d\phi'\\&=-\uvec{x}\frac{\mu_0Ib^2}{4R^2}\sin\theta\end{aligned}}

\begin{itemize}
\tightlist
\item
  At an arbitrary field point, the vector magnetic potential is then
\end{itemize}

\ELdisp{\vect{A}\ELbr =\uvec{\phi}\frac{\mu_0Ib^2}{4R^2}\sin\theta\quad\ELbr \ELbr \Longrightarrow\quad\ELbr \vect{A}\ELbr =\uvec{\phi}\frac{\mu_0(I\pi b^2)}{4\pi R^2}\sin\theta}

\begin{itemize}
\tightlist
\item
  The magnetic dipole moment vector
\end{itemize}

\ELdisp{\vect{m}\ELbr =\uvec{z}I\pi b^2\ELbr =\uvec{z}IS\ELbr =\uvec{z}m}

\begin{ELimportant}{}

\ELdisp{\vect{A}\ELbr =\frac{\mu_0\vect{m}\times\uvec{R}}{4\pi R^2}}

\end{ELimportant}

\ELdisp{\vect{B}\ELbr =\nabla\times\vect{A}\quad\ELbr \ELbr \Longrightarrow\quad\ELbr \vect{B}\ELbr =\frac{\mu_0Ib^2}{4R^3}\left(\uvec{R}2\cos\theta+\uvec{\theta}\sin\theta\right)}

\begin{ELimportant}{}

\ELdisp{\vect{B}\ELbr =\frac{\mu_0m}{4\pi R^3}\left(\uvec{R}2\cos\theta+\uvec{\theta}\sin\theta\right)}

\end{ELimportant}

\ELfigure{m05-dipoles}{}

\end{ELsolution}

\end{ELexample}

\ELhold

\ELsection{The Static Magnetic Field in Material Media}{The Static Magnetic Field in Material Media}

\begin{ELunit}

\ELfigure{m05-map-media}{}

\begin{itemize}
\tightlist
\item
  All materials consist of atoms with a positively charged nucleus and a number of negatively charged electrons orbiting around it.
\item
  The motion of the electrons produces circulating currents and hence microscopic magnetic dipoles.
\item
  In the absence of an external magnetic field, the magnetic dipoles of the atoms of most materials (except permanent magnets) have random orientations, and no net magnetic moment results. When an external magnetic field is applied, the magnetic moments of the orbiting electrons align.
\item
  How do different materials behave when exposed to a magnetic field, and how can this behaviour be analysed?
\end{itemize}

\end{ELunit}

\ELhold

\ELsubsection{Magnetization and Equivalent Current Densities}{Magnetization and Equivalent Current Densities}

\begin{ELunit}

\ELfigure{m05-magnetization}{}

\begin{itemize}
\tightlist
\item
  The magnetization vector (and the polarization vector in electrostatics)
\end{itemize}

\ELdisp{\vect{M}\ELbr =\lim_{\Delta v\to0}\frac{\sum_{k=1}^{n\Delta v}\vect{m}_k}{\Delta v}\quad\ELbr (\mathrm{A/m})\qquad\ELbr \qquad\ELbr \vect{P}\ELbr =\lim_{\Delta v\to0}\frac{\sum_{k=1}^{n\Delta v}\vect{p}_k}{\Delta v}\quad\ELbr (\mathrm{C/m^2})}

\begin{itemize}
\tightlist
\item
  The vector magnetic potential of a magnetized material (and the electric potential of a polarized material)
\end{itemize}

\ELdisp{\vect{A}\ELbr =\frac{\mu_0}{4\pi}\int_{V'}\frac{\nabla'\times\vect{M}}{R}\,dv'+\frac{\mu_0}{4\pi}\oint_{S'}\frac{\vect{M}\times\uvec{n}'}{R}\,ds'}
\ELdisp{V\ELbr =\frac{1}{4\pi\epsilon_0}\oint_{S'}\frac{\vect{P}\cdot\uvec{n}'}{R}\,ds'+\frac{1}{4\pi\epsilon_0}\int_{V'}\frac{(-\nabla'\cdot\vect{P})}{R}\,dv'}

\end{ELunit}

\begin{ELjoin}

\begin{itemize}
\tightlist
\item
  The equivalent magnetization surface and volume current densities
\end{itemize}

\begin{ELimportant}{}

\ELdisp{\vect{J}_{ms}\ELbr =\vect{M}\times\uvec{n},\qquad\ELbr \vect{J}_m\ELbr =\nabla\times\vect{M}}

\end{ELimportant}

\end{ELjoin}

\begin{ELunit}

\begin{itemize}
\tightlist
\item
  The equivalent polarization surface and volume charge densities
\end{itemize}

\ELdisp{\rho_{ps}\ELbr =\vect{P}\cdot\uvec{n},\qquad\ELbr \rho_p\ELbr =-\nabla\cdot\vect{P}}

\end{ELunit}

\begin{ELjoin}

\begin{itemize}
\tightlist
\item
  The problem of finding the magnetic flux density \ELim{\vect{B}} caused by a given volume density of magnetic dipole moment \ELim{\vect{M}} becomes one of finding the equivalent magnetization current densities, then \ELim{\vect{A}}, and then \ELim{\vect{B}}.
\end{itemize}

\ELnextnum{5-8}

\begin{ELexample}{}

Determine the magnetic flux density on the axis of a magnetized cylinder. The cylinder has a radius \ELim{b}, a length \ELim{L} and a magnetization vector \ELim{\vect{M}=\uvec{z}M_0}.

\ELfigure{m05-ex8}{}

\begin{ELsolution}

\ELdisp{\vect{J}_m\ELbr =\nabla\times\vect{M}\ELbr =0}

\begin{itemize}
\tightlist
\item
  The equivalent surface current density on the side wall
\end{itemize}

\ELdisp{\vect{J}_{ms}\ELbr =\vect{M}\times\uvec{n}\ELbr =\uvec{z}M_0\times\uvec{r}\ELbr =\uvec{\phi}M_0}

\begin{itemize}
\tightlist
\item
  The equivalent surface current density on the top and bottom faces is zero.
\end{itemize}

\ELdisp{\left.\begin{aligned}&\vect{B}=\frac{\mu_0}{4\pi}\int_{S'}\frac{\vect{J}_{ms}\times(\vect{R}-\vect{R}')\,ds'}{\lvert\vect{R}-\vect{R}'\rvert^3}\\&ds'=b\,d\phi'\,dz'\\&\vect{R}=z\uvec{z}\\&\vect{R}'=b\uvec{r'}+z'\uvec{z}\\&\vect{R}-\vect{R}'=(z-z')\uvec{z}-b\uvec{r'}\\&\begin{aligned}\vect{J}_{ms}\times(\vect{R}-\vect{R}')&=M_0\uvec{\phi'}\times\left((z-z')\uvec{z}-b\uvec{r'}\right)\\&=M_0(z-z')\uvec{r'}+bM_0\uvec{z}\end{aligned}\end{aligned}\right\}}

\begin{itemize}
\tightlist
\item
  It is evident from the geometry that the \ELim{r}-component of the magnetic flux density is zero.
\end{itemize}

\ELdisp{\begin{aligned}\vect{B}=\int d\vect{B}&=\uvec{z}\int_0^L\frac{\mu_0M_0b^2\,dz'}{2\left[(z-z')^2+b^2\right]^{3/2}}\\&=\uvec{z}\frac{\mu_0M_0}{2}\left[\frac{z}{\sqrt{z^2+b^2}}-\frac{z-L}{\sqrt{(z-L)^2+b^2}}\right]\end{aligned}}

\end{ELsolution}

\end{ELexample}

\end{ELjoin}

\ELhold

\ELsection{Magnetic Field Intensity and Relative Permeability}{Magnetic Field Intensity and Relative Permeability}

\begin{ELunit}

\ELfigure{m05-map-hmu}{}

\end{ELunit}

\ELhold

\ELsubsection{Magnetic Field Intensity}{Magnetic Field Intensity}

\begin{ELunit}

\begin{itemize}
\tightlist
\item
  Since an external magnetic field aligns the internal dipole moments and induces a magnetic moment in a magnetic material, we expect the resulting magnetic flux density in the presence of the magnetic material to differ from its value in free space.
\item
  Taking the equivalent volume current density in the material into account, we have
\end{itemize}

\ELdisp{\frac{1}{\mu_0}\nabla\times\vect{B}\ELbr =\vect{J}+\vect{J}_m\ELbr =\vect{J}+\nabla\times\vect{M}\quad\ELbr \ELbr \Longrightarrow\quad\ELbr \nabla\times\left(\frac{\vect{B}}{\mu_0}-\vect{M}\right)\ELbr =\vect{J}}

\end{ELunit}

\begin{ELjoin}

\begin{itemize}
\tightlist
\item
  We define the magnetic field intensity as
\end{itemize}

\begin{ELimportant}{}

\ELdisp{\vect{H}\ELbr =\frac{\vect{B}}{\mu_0}-\vect{M}\qquad\ELbr (\mathrm{A/m})}

\end{ELimportant}

\end{ELjoin}

\begin{ELjoin}

\begin{itemize}
\tightlist
\item
  The vector \ELim{\vect{H}} lets us write the curl equation relating the magnetic field and the distribution of free currents in any medium without using \ELim{\vect{M}} or \ELim{\vect{J}_m}
\end{itemize}

\begin{ELimportant}{}

\ELdisp{\nabla\times\vect{H}\ELbr =\vect{J}\qquad\ELbr (\mathrm{A/m^2})}

\end{ELimportant}

\end{ELjoin}

\begin{ELjoin}

\begin{itemize}
\tightlist
\item
  \ELim{\vect{J}}: the density of the free current
\item
  The integral form
\end{itemize}

\ELdisp{\int_S(\nabla\times\vect{H})\cdot d\vect{s}\ELbr =\int_S\vect{J}\cdot d\vect{s}\quad\ELbr \ELbr \Longrightarrow\quad\ELbr \boxed{\oint_C\vect{H}\cdot\dif\vect{\ell}=I\qquad(\mathrm{A})}}

\begin{itemize}
\tightlist
\item
  This last relation is another form of Ampère's circuital law.
\item
  The two fundamental governing equations for magnetostatics in any medium
\end{itemize}

\begin{ELimportant}{}

\ELdisp{\left\{\begin{aligned}&\nabla\cdot\vect{B}=0\\&\nabla\times\vect{H}=\vect{J}\end{aligned}\right.}

\end{ELimportant}

\end{ELjoin}

\ELhold

\ELsubsection{Relative Permeability}{Relative Permeability}

\begin{ELjoin}

\begin{itemize}
\tightlist
\item
  The relation between the magnetization vector and the magnetic field intensity
\end{itemize}

\ELdisp{\vect{M}\ELbr =\chi_m\vect{H}}

\begin{itemize}
\tightlist
\item
  \ELim{\chi_m}: the magnetic susceptibility
\item
  The relation between the magnetic field intensity and the magnetic flux density
\end{itemize}

\begin{ELimportant}{}

\ELdisp{\begin{aligned}\vect{B}&=\mu_0(1+\chi_m)\vect{H}\\&=\mu_0\mu_r\vect{H}=\mu\vect{H}\qquad(\mathrm{Wb/m^2})\end{aligned}}

\end{ELimportant}

\end{ELjoin}

\begin{ELunit}

\ELdisp{\mu_r\ELbr =1+\chi_m\ELbr =\frac{\mu}{\mu_0}}

\begin{itemize}
\tightlist
\item
  \ELim{\mu_r}: a dimensionless quantity called the relative permeability of the medium
\item
  The relations of electrostatics and of magnetostatics are dual to each other
\end{itemize}

\end{ELunit}

\Needspace{19\baselineskip}
\begin{ELltr}
\begin{longtable}[c]{cc}
\toprule\noalign{}
\ELtablehead Electrostatics & Magnetostatics \\\noalign{\vskip 4pt}
\midrule\noalign{}
\endhead
\bottomrule\noalign{}
\endlastfoot
\ELim{\vect{E}} & \ELim{\vect{B}} \\\noalign{\vskip 4pt}
\ELim{\vect{D}} & \ELim{\vect{H}} \\\noalign{\vskip 4pt}
\ELim{\epsilon} & \ELim{\dfrac1\mu} \\\noalign{\vskip 4pt}
\ELim{\vect{P}} & \ELim{-\vect{M}} \\\noalign{\vskip 4pt}
\ELim{\rho} & \ELim{\vect{J}} \\\noalign{\vskip 4pt}
\ELim{V} & \ELim{\vect{A}} \\\noalign{\vskip 4pt}
\ELim{\cdot} & \ELim{\times} \\\noalign{\vskip 4pt}
\ELim{\times} & \ELim{\cdot} \\\noalign{\vskip 4pt}
\end{longtable}\end{ELltr}


\ELhold

\ELsubsection{Behaviour of Magnetic Materials}{Behaviour of Magnetic Materials}

\begin{ELunit}

\begin{itemize}
\tightlist
\item
  Magnetic materials can be roughly classified into three main groups according to their relative permeability

  \begin{itemize}
  \tightlist
  \item
    \textbf{Diamagnetic materials}: \ELim{\mu_r<1} (\ELim{\chi_m} is a very small negative number)

    \begin{itemize}
    \tightlist
    \item
      In these materials the net magnetic moment is zero in the absence of an applied external magnetic field. \ELim{\ELto} copper, lead, mercury, silver and gold \ELim{\ELto} \ELim{\chi_m\approx-10^{-5}}
    \end{itemize}
  \item
    \textbf{Paramagnetic materials}: \ELim{\mu_r>1} (\ELim{\chi_m} is a very small positive number)

    \begin{itemize}
    \tightlist
    \item
      In these materials there is a net average magnetic moment in the absence of an applied external magnetic field. \ELim{\ELto} aluminium, magnesium and tungsten \ELim{\ELto} \ELim{\chi_m\approx10^{-5}}
    \end{itemize}
  \item
    \textbf{Ferromagnetic materials}: \ELim{\mu_r\gg1} (\ELim{\chi_m} is a large positive number)

    \begin{itemize}
    \tightlist
    \item
      In these materials the magnetic moment is many times larger than in paramagnetic materials. \ELim{\ELto} cobalt, nickel and iron
    \item
      The relation between \ELim{\vect{B}} and \ELim{\vect{H}} in a ferromagnetic material is nonlinear.
    \end{itemize}
  \end{itemize}
\end{itemize}

\ELfigure{m05-hysteresis}{The hysteresis phenomenon}

\end{ELunit}

\ELhold

\ELsection{Magnetic Circuits}{Magnetic Circuits}

\begin{ELunit}

\ELfigure{m05-map-circ}{}

\begin{itemize}
\tightlist
\item
  As with electric circuits, it is sometimes necessary to calculate the magnetic flux and the magnetic field intensity produced in a magnetic circuit by current-carrying windings.
\item
  To analyse magnetic circuits we have
\end{itemize}

\ELdisp{\nabla\cdot\vect{B}\ELbr =0,\qquad\ELbr \nabla\times\vect{H}\ELbr =\vect{J}}
\ELdisp{\oint_C\vect{H}\cdot\dif\vect{\ell}\ELbr =NI\ELbr =\mathcal{V}_m}

\end{ELunit}

\begin{ELjoin}

\begin{itemize}
\tightlist
\item
  \ELim{\mathcal{V}_m}: the magnetomotive force, in \ELim{\mathrm{A}}. It is the analogue of the electromotive force in electric circuits.
\end{itemize}

\ELnextnum{5-9}

\begin{ELexample}{}

Assume that \ELim{N} turns of wire are wound around a toroidal core of a magnetic material with permeability \ELim{\mu}. The core has a mean radius \ELim{r_0}, a circular cross-section of radius \ELim{a} and a narrow air gap of length \ELim{\ell_g}. A steady current \ELim{I_0} flows in the wire. Determine

\begin{itemize}
\item
  \begin{enumerate}
  \def\labelenumi{(\alph{enumi})}
  \tightlist
  \item
    the magnetic flux density in the magnetic core,
  \end{enumerate}
\item
  \begin{enumerate}
  \def\labelenumi{(\alph{enumi})}
  \setcounter{enumi}{1}
  \tightlist
  \item
    the magnetic field intensity in the magnetic core,
  \end{enumerate}
\item
  \begin{enumerate}
  \def\labelenumi{(\alph{enumi})}
  \setcounter{enumi}{2}
  \tightlist
  \item
    the magnetic field intensity in the air gap.
  \end{enumerate}
\end{itemize}

\ELfigure{m05-ex9}{}

\begin{ELsolution}

\begin{itemize}
\tightlist
\item
  \textbf{First assumption}: the leakage flux is neglected
\end{itemize}

\ELfigure{m05-ex9-gauss}{}

\ELdisp{\nabla\cdot\vect{B}\ELbr =0\quad\ELbr \ELbr \Longrightarrow\quad\ELbr \oint_S\vect{B}\cdot d\vect{s}\ELbr =0}

\begin{itemize}
\tightlist
\item
  The flux in the core and in the air gap is the same.
\item
  \textbf{Second assumption}: the fringing of the flux in the air gap is neglected.

  \begin{itemize}
  \tightlist
  \item
    The flux density in the core and in the air gap is the same.
  \end{itemize}
\end{itemize}

\ELdisp{\vect{B}_f\ELbr =\vect{B}_g\ELbr =\uvec{\phi}B_f}
\ELdisp{\vect{H}_f\ELbr =\uvec{\phi}\frac{B_f}{\mu},\qquad\ELbr \vect{H}_g\ELbr =\uvec{\phi}\frac{B_f}{\mu_0}}
\ELdisp{\oint_C\vect{H}\cdot\dif\vect{\ell}\ELbr =NI_0\quad\ELbr \ELbr \Longrightarrow\quad\ELbr \frac{B_f}{\mu}(2\pi r_0-\ell_g)+\frac{B_f}{\mu_0}\ell_g\ELbr =NI_0}

\begin{itemize}
\item
  \begin{enumerate}
  \def\labelenumi{(\alph{enumi})}
  \tightlist
  \item
  \end{enumerate}
\end{itemize}

\ELdisp{\vect{B}_f\ELbr =\uvec{\phi}\frac{\mu_0\mu NI_0}{\mu_0(2\pi r_0-\ell_g)+\mu\ell_g}}

\begin{itemize}
\item
  \begin{enumerate}
  \def\labelenumi{(\alph{enumi})}
  \setcounter{enumi}{1}
  \tightlist
  \item
  \end{enumerate}
\end{itemize}

\ELdisp{\vect{H}_f\ELbr =\uvec{\phi}\frac{\mu_0NI_0}{\mu_0(2\pi r_0-\ell_g)+\mu\ell_g}}

\begin{itemize}
\item
  \begin{enumerate}
  \def\labelenumi{(\alph{enumi})}
  \setcounter{enumi}{2}
  \tightlist
  \item
  \end{enumerate}
\end{itemize}

\ELdisp{\vect{H}_g\ELbr =\uvec{\phi}\frac{\mu NI_0}{\mu_0(2\pi r_0-\ell_g)+\mu\ell_g}}

\end{ELsolution}

\end{ELexample}

\end{ELjoin}

\begin{ELjoin}

\begin{itemize}
\tightlist
\item
  If the radius of the cross-section of the core is much smaller than the mean radius of the toroid, the magnetic flux density in the core is approximately constant, and
\end{itemize}

\ELdisp{\Phi\ELbr \cong BS\quad\ELbr \ELbr \Longrightarrow\quad\ELbr \Phi\ELbr =\frac{NI_0}{(2\pi r_0-\ell_g)/\mu S+\ell_g/\mu_0S}\quad\ELbr \ELbr \Longrightarrow\quad\ELbr \Phi\ELbr =\frac{\mathcal{V}_m}{\mathcal{R}_f+\mathcal{R}_g}}

\begin{ELimportant}{}

\ELdisp{\mathcal{R}_f\ELbr =\frac{2\pi r_0-\ell_g}{\mu S}\ELbr =\frac{\ell_f}{\mu S},\qquad\ELbr \mathcal{R}_g\ELbr =\frac{\ell_g}{\mu_0S}}

\end{ELimportant}

\end{ELjoin}

\begin{ELunit}

\begin{itemize}
\item
  \ELim{\mathcal{R}_f} and \ELim{\mathcal{R}_g}: reluctances, in \ELim{1/\mathrm{H}}
\item
  This magnetic circuit is analogous to an electric circuit as follows
\end{itemize}

\ELfigure{m05-circuits}{}

\ELdisp{\Phi\ELbr =\frac{\mathcal{V}_m}{\mathcal{R}_f+\mathcal{R}_g},\qquad\ELbr  I\ELbr =\frac{\mathcal{V}}{R_f+R_g}}

\end{ELunit}

\Needspace{12\baselineskip}
\begin{ELltr}
\begin{longtable}[c]{ll}
\toprule\noalign{}
\ELtablehead Magnetic Circuits & Electric Circuits \\\noalign{\vskip 4pt}
\midrule\noalign{}
\endhead
\bottomrule\noalign{}
\endlastfoot
mmf, \ELim{\mathcal{V}_m\,(=NI)} & emf, \ELim{\mathcal{V}} \\\noalign{\vskip 4pt}
magnetic flux, \ELim{\Phi} & electric current, \ELim{I} \\\noalign{\vskip 4pt}
reluctance, \ELim{\mathcal{R}} & resistance, \ELim{R} \\\noalign{\vskip 4pt}
permeability, \ELim{\mu} & conductivity, \ELim{\sigma} \\\noalign{\vskip 4pt}
\end{longtable}\end{ELltr}


\begin{ELjoin}

\begin{itemize}
\tightlist
\item
  In analogy to Kirchhoff's voltage law in electric circuits, in magnetic circuits we have
\end{itemize}

\begin{ELimportant}{}

\ELdisp{\sum_jN_jI_j\ELbr =\sum_k\mathcal{R}_k\Phi_k}

\end{ELimportant}

\end{ELjoin}

\begin{ELjoin}

\begin{itemize}
\tightlist
\item
  So around any closed path in a magnetic circuit, the algebraic sum of the ampere-turns equals the algebraic sum of the products of the reluctances and the fluxes.
\item
  In analogy to Kirchhoff's current law in electric circuits, in magnetic circuits we have
\end{itemize}

\begin{ELimportant}{}

\ELdisp{\sum_j\Phi_j\ELbr =0}

\end{ELimportant}

\end{ELjoin}

\begin{ELunit}

\begin{itemize}
\tightlist
\item
  which states that the algebraic sum of all the magnetic fluxes flowing out of a junction in a magnetic circuit is zero.
\item
  Despite this simple analogy, an exact analysis of magnetic circuits is very difficult, for the following reasons

  \begin{itemize}
  \tightlist
  \item
    the leakage fluxes,
  \item
    the fringing effects, which make the magnetic flux lines spread out and bulge in the air gap,
  \item
    the dependence of the permeability of ferromagnetic materials on the magnetic field intensity, that is, the nonlinear relation between \ELim{\vect{B}} and \ELim{\vect{H}}.
  \end{itemize}
\end{itemize}

\end{ELunit}

\ELnextnum{5-10}

\begin{ELexample}{}

Consider the magnetic circuit shown below. Steady currents \ELim{I_1} and \ELim{I_2} flow in windings of \ELim{N_1} and \ELim{N_2} turns respectively. The core has a cross-sectional area \ELim{S_c} and a permeability \ELim{\mu}. Determine the magnetic flux in the middle leg.

\ELfigure{m05-ex10}{}

\begin{ELsolution}

\begin{itemize}
\tightlist
\item
  The equivalent circuit
\end{itemize}

\ELfigure{m05-ex10-circuit}{}

\ELdisp{\mathcal{R}_1\ELbr =\frac{\ell_1}{\mu S_c},\qquad\ELbr \mathcal{R}_2\ELbr =\frac{\ell_2}{\mu S_c},\qquad\ELbr \mathcal{R}_3\ELbr =\frac{\ell_3}{\mu S_c}}
\ELdisp{\begin{aligned}&\mathit{Loop}~1:&N_1I_1&=(\mathcal{R}_1+\mathcal{R}_3)\Phi_1+\mathcal{R}_1\Phi_2\\&\mathit{Loop}~2:&N_1I_1-N_2I_2&=\mathcal{R}_1\Phi_1+(\mathcal{R}_1+\mathcal{R}_2)\Phi_2\end{aligned}}
\ELdisp{\Phi_1\ELbr =\frac{\mathcal{R}_2N_1I_1+\mathcal{R}_1N_2I_2}{\mathcal{R}_1\mathcal{R}_2+\mathcal{R}_1\mathcal{R}_3+\mathcal{R}_2\mathcal{R}_3}}

\end{ELsolution}

\end{ELexample}

\ELhold

\ELsection{Boundary Conditions for Magnetostatic Fields}{Boundary Conditions for Magnetostatic Fields}

\begin{ELunit}

\ELfigure{m05-map-bc}{}

\begin{itemize}
\tightlist
\item
  The boundary condition for the normal components
\end{itemize}

\ELdisp{\nabla\cdot\vect{B}\ELbr =0\quad\ELbr \ELbr \Longrightarrow\quad\ELbr \boxed{B_{1n}=B_{2n}\qquad(\mathrm{T})}}

\end{ELunit}

\begin{ELjoin}

\begin{itemize}
\tightlist
\item
  In linear media
\end{itemize}

\begin{ELimportant}{}

\ELdisp{\mu_1H_{1n}\ELbr =\mu_2H_{2n}}

\end{ELimportant}

\end{ELjoin}

\begin{ELunit}

\begin{itemize}
\tightlist
\item
  The boundary condition for the tangential components
\end{itemize}

\ELfigure{m05-bc}{}

\ELdisp{\oint_C\vect{H}\cdot\dif\vect{\ell}\ELbr =I}
\ELdisp{\oint_{abcda}\vect{H}\cdot\dif\vect{\ell}\ELbr =\vect{H}_1\cdot\Delta\vect{w}+\vect{H}_2\cdot(-\Delta\vect{w})\ELbr =J_{sn}\Delta w}
\ELdisp{H_{1t}-H_{2t}\ELbr =J_{sn}\qquad\ELbr (\mathrm{A/m})}

\end{ELunit}

\begin{ELimportant}{}

\ELdisp{\uvec{n2}\times(\vect{H}_1-\vect{H}_2)\ELbr =\vect{J}_s\qquad\ELbr (\mathrm{A/m})}

\end{ELimportant}
